\documentclass[10pt,a4paper]{amsart}
\usepackage[margin=19mm]{geometry}
\usepackage{amsmath,amssymb,mathtools}
\usepackage[scr=rsfs,cal=euler]{mathalfa}
\usepackage{xfrac}
\usepackage[unicode,hidelinks,bookmarks=false]{hyperref}
\newcommand{\ie}{i.e.\ }
\newcommand{\cf}{cf.\ }
\newcommand{\resp}{respectively\ }
\newcommand{\Subsection}{\subsection}
\providecommand{\nobreakdash}{\nobreak\hbox{-}}
\setlength{\emergencystretch}{2em}
\multlinegap0pt
\usepackage{bm}
\usepackage{bbm}
\usepackage{dsfont}
\usepackage{xspace}
\usepackage{cancel}
\usepackage{mathtools}
\usepackage{amsthm}
\makeatletter
\@ifundefined{theorem}{\newtheorem{theorem}{Theorem}}{}
\@ifundefined{lemma}{\newtheorem{lemma}[theorem]{Lemma}}{}
\makeatother
\DeclareMathOperator{\dist}{dist}
\newcommand{\field}{\mathbb{C}}
\newcommand{\identity}{\mathbf{1}}
\title[The Multiplicative Kowalski-Slodkowski Theorem for Hermitian]{The Multiplicative Kowalski-Slodkowski Theorem for Hermitian Algebras}
\author{Rudi Brits, Muhammad Hassen, Cheick Toure}
\date{Source: arXiv:2509.03663v2 (CC BY 4.0)}
\begin{document}
\maketitle

\noindent\emph{Source and licence.} The Multiplicative Kowalski-Slodkowski Theorem for Hermitian Algebras. Version arXiv:2509.03663v2 (CC BY 4.0), distributed under the Creative Commons Attribution 4.0 International licence, as recorded in the arXiv licence metadata for this exact version and shown on the abstract page https://arxiv.org/abs/2509.03663v2. Full rights notice: licenses/arxiv-2509.03663-CC-BY-4.0.txt.\par\smallskip

\noindent\textit{Editorial scope.} Counted unit: the Lemma whose source label is \texttt{no\_zero} in the article (source0.tex lines 241--244, source label \texttt{no\_zero}), delivered together with the article's own complete original proof of it (source0.tex lines 245--264, one \texttt{proof} environment of 20 lines). No mathematics is written, repaired or re-derived: statement and proof are the article's own text line for line. Required context retained uncounted: lem:1 (lines 169--171); lem:2 (lines 181--183); lem:3 (lines 195--202); lem:5 (lines 227--229). Every definition, display and label of those lines is retained unchanged. Omitted from this package and not redistributed: 1--168, 172--180, 184--194, 203--226, 230--240, 265--355. Nothing inside a retained excerpt is omitted or altered: every retained body is delivered exactly as the article's lines are written, so no omission receipt of any kind accompanies this package. Citations: the retained excerpts carry no citation command at all, so no bibliography entry of the article is transcribed and no reference is redistributed. Reference adaptation: none was needed and none was made; every reference inside the delivered unit resolves inside it, so no unresolved reference or citation is produced. Printed slips kept literal: no printed word, symbol, hypothesis or number of the counted statement or of its proof was altered. Layout and class: the article's own class is \texttt{amsart} with packages including amssymb, amsmath, amsthm, mathtools, hyperref, titlesec, amsaddr; the wrapper is the lane's pinned \texttt{amsart} editorial wrapper at 10pt on A4, which restates the theorem-like environments the retained text needs and adds the article's own macro definitions that the retained text needs (\texttt{\textbackslash dist}, \texttt{\textbackslash field}, \texttt{\textbackslash identity}), with extra wrapper packages (bm, bbm, dsfont, xspace, cancel, mathtools). The delivered numbering is the unit's own. Assets: no retained span contains a figure environment, a graphics-inclusion command, a PDF-page inclusion, a table or a TikZ picture; no image file is copied into this package, the package declares no asset dependency and no image file is redistributed with it. Licence: version arXiv:2509.03663v2 of this article is distributed under the Creative Commons Attribution 4.0 International licence, as recorded in the arXiv licence metadata for this exact version and shown on the abstract page https://arxiv.org/abs/2509.03663v2. A full rights notice is delivered at licenses/arxiv-2509.03663-CC-BY-4.0.txt. Characters: every retained excerpt is pure ASCII, so the delivered engine, XeLaTeX with its default Latin Modern text font, reports no missing character.
	\begin{lemma}\label{lem:1}
		Let $\left(a_n\right)$ and $\left(b_n\right)$ be sequences in $S_A$ such that $a_n^2 < b_n^2$ for all $n\in \mathbb{N}$. If $\left(v_n\right)$ is a sequence in $A$ such that $s(b_nv_n)\to 0$, then $s(a_nv_n)\to 0$.
	\end{lemma}

	\begin{lemma}\label{lem:2}
			Let $a\in S_A$, and let $(v_n)$ be a sequence in $A$ such that $s\left(\sqrt{a^2+\frac{1}{n^2}\mathbf 1}\, v_n\right)\to 0$. Then $s(av_n)\to 0$.
	\end{lemma}

	\begin{lemma}\label{lem:3}
		Let $0\leq a$. Then 
		\begin{itemize}
		\item[\rm (a)] $\lim\limits_{n\to\infty}s\left(\sqrt{a+\frac{1}{n^2}\mathbf 1}\,e^{-n\sqrt{a+\frac{1}{n^2}\mathbf 1}}\right)=0,$
		\item[\rm (b)]
		$\lim\limits_{n\to\infty}s \left(\sqrt{a+\frac{1}{n^2}\mathbf 1}\left(\identity+in\sqrt{a+\frac{1}{n^2}\mathbf 1}\right)^{-1}\right)=0.$
		\end{itemize}
	\end{lemma}

 	\begin{lemma}\label{lem:5}
 		Let $\phi$ be a continuous spectrally multiplicative function on $A$ with $\phi(\identity)=1$. If $\alpha\in\field$, and $x\in A$ satisfies $\psi_\phi(x)=0$, then $\phi(\alpha\identity+x)=c_\alpha \alpha$ for some $c_\alpha\in [0,1]$.
 	\end{lemma}

  	\begin{lemma}\label{no_zero}
  		Let  $\phi$ be a 
  		continuous   spectrally multiplicative  function $A$ with $\phi(\identity)=1$. If $\alpha \in \mathbb{C},$ and  $x\in A$  satisfies $\psi_\phi(x)=0 $, then $\phi(\alpha\mathbf{1}+ x)\in\{ 0,\alpha  \}.$
  	\end{lemma}

  	\begin{proof}
  		For each $n\in\mathbb N$ let $V_n:=\left(\mathbf{1}+in\sqrt{R^2+I^2+\frac{1}{n^2}\identity} \right)^{-1}.$
  		Again using Lemma~\ref{lem:1}, Lemma~\ref{lem:2} and Lemma~ \ref{lem:3}, we have that $$\lim_ns\left(\sqrt{R^2+I^2+\frac{1}{n^2}\identity}\, V_n\right)={0} \implies \lim_ns\left( R  V_n\right)=0.$$
  		Similarly  $\lim_ns\left( I  V_n\right)=0 $. Observe that each $V_n$ belongs to $ G_{\mathbf 1}(A)$, whence it follows that $\phi(V_n) = \psi_\phi(V_n)=1/2-i/2 $. Let $\alpha \neq 0 $. From Lemma~\ref{lem:5}, we have that $\phi(\alpha\mathbf{1}+x)=c_\alpha \alpha $, with $c_\alpha \in [0,1]$. To obtain the result we have to show that $c_\alpha \in \{0,1\}$: For the sake of a contradiction assume that $ 0< c_\alpha <1$. If we set $Z_n:=\frac{1}{\alpha}xV_n=\frac{1}{\alpha}(R+iI) V_n$, then
  		
  		\begin{equation}\label{3}
  			c_\alpha(1/2-i/2)=\frac{1}{\alpha}\phi(\alpha\mathbf{1}+x)\phi(V_n)\in \sigma\left(V_n+Z_n\right).
  		\end{equation}
  		
  		The first paragraph of the proof shows that $\lim_n s\left(Z_n\right)=0$, and \eqref{3} shows that 
  		$c_\alpha(1/2-i/2)\mathbf{1}-V_n-Z_n  \notin G(A) $. From the definition of $V_n$ together with the spectral mapping theorem, and the fact that $0<\sqrt{R^2+I^2+\frac{1}{n^2}\identity}$ we have that $\sigma\left(V_n\right) \subseteq C_r $, where $C_r$ is the circle in $\mathbb C$ with center $\frac{1}{2}$ and radius $\frac{1}{2}$. Thus $$c_\alpha(1/2-i/2)\mathbf{1} -V_n-Z_n \notin G(A) \mbox{ and } c_\alpha(1/2-i/2)\mathbf{1} -V_n \in G(A) ,$$ which together implies that 
  		\begin{equation}\label{4}
  		 \mathbf{1}-Z_n\left( c_\alpha(1/2-i/2)\mathbf{1} -V_n\right)^{-1} \notin G(A).
  		\end{equation}
  		Since $V_n$ is normal we obtain the estimate
  		\begin{align*}
  		s\left(\left( c_\alpha(1/2-i/2)\mathbf{1} -V_n\right)^{-1}\right)&=\rho\left( \left( c_\alpha(1/2-i/2)\mathbf{1} -V_n\right)^{-1} \right)\\& \leq \frac{1}{\dist(C_r,\{ c_\alpha(1/2-i/2) \} ) }
  		\end{align*}
  		from which it follows that $\lim_ns\left(  Z_n\left( c_\alpha(1/2-i/2)\mathbf{1} -V_n\right)^{-1}\right)=0$, contradicting \eqref{4}. Subsequently $c_\alpha \in \{0,1\}$, and $\phi(\alpha\mathbf{1}+x)\in \{0, \alpha \}$ follows as advertised. 
  	\end{proof}

\end{document}
